%%
%% Copyright 2019-2024 Elsevier Ltd
%%
%% This file is part of the 'CAS Bundle'.
%% --------------------------------------
%%
%% It may be distributed under the conditions of the LaTeX Project Public
%% License, either version 1.3c of this license or (at your option) any
%% later version.  The latest version of this license is in
%%    http://www.latex-project.org/lppl.txt
%% and version 1.3c or later is part of all distributions of LaTeX
%% version 1999/12/01 or later.
%%
%% The list of all files belonging to the 'CAS Bundle' is
%% given in the file `manifest.txt'.
%%
%% Template article for cas-dc documentclass for
%% double column output.

\documentclass[a4paper,fleqn]{cas-dc}

% If the frontmatter runs over more than one page
% use the longmktitle option.

%\documentclass[a4paper,fleqn,longmktitle]{cas-dc}

%\usepackage[numbers]{natbib}
%\usepackage[authoryear]{natbib}
\usepackage[authoryear,longnamesfirst]{natbib}

%%%Author macros
\def\tsc#1{\csdef{#1}{\textsc{\lowercase{#1}}\xspace}}
\tsc{WGM}
\tsc{QE}
%%%

% Uncomment and use as if needed
%\newtheorem{theorem}{Theorem}
%\newtheorem{lemma}[theorem]{Lemma}
%\newdefinition{rmk}{Remark}
%\newproof{pf}{Proof}
%\newproof{pot}{Proof of Theorem \ref{thm}}

\begin{document}
\let\WriteBookmarks\relax
\def\floatpagepagefraction{1}
\def\textpagefraction{.001}

% Short title
\shorttitle{Physics-informed LSTM and BiGRU for satellite trajectory prediction}

% Short author
\shortauthors{aaa \emph{et~al.}}

% Main title of the paper
\title [mode = title]{Physics-informed LSTM and BiGRU-based approach for \\
maneuverable satellite cluster trajectory prediction and intention recognition}

% Title footnote mark
\tnotemark[1]

% Title footnote 1.
\tnotetext[1]{Submitted to \emph{Aerospace Science and Technology}.}

% First author
\author[1]{aaa}[
    orcid=0000-0000-0000-0000,
]

% Corresponding author indication
\cormark[1]

% Footnote of the first author
\fnmark[1]

% Email id of the first author
\ead{aaa@example.edu}

% Credit authorship
\credit{Conceptualization, Methodology, Software, Writing -- original draft}

% Address/affiliation
\affiliation[1]{organization={School of Astronautics, Example University},
            addressline={No.~100 University Road},
            city={Beijing},
            postcode={100000},
            state={},
            country={China}}

% Second author
\author[1]{bbb}[
    orcid=0000-0000-0000-0001,
]

\fnmark[1]

\ead{bbb@example.edu}

\credit{Formal analysis, Validation, Writing -- review \& editing}

% Third author
\author[2]{ccc}[
    orcid=0000-0000-0000-0002,
]

\fnmark[1]

\ead{ccc@example.edu}

\credit{Data curation, Investigation, Visualization}

\affiliation[2]{organization={Institute of Aerospace Engineering, Example Academy},
            addressline={No.~200 Science Avenue},
            city={Shanghai},
            postcode={200000},
            state={},
            country={China}}

% Corresponding author text
\cortext[1]{Corresponding author.}

% Footnote text
\fntext[1]{These authors contributed equally to this work.}

% For a title note without a number/mark
%\nonumnote{}

% Here goes the abstract
\begin{abstract}
This paper proposes a novel approach combining physics-informed long short-term
memory (LSTM) networks and bidirectional gated recurrent units (BiGRU) for
maneuverable satellite cluster trajectory prediction and intention recognition.
The physics-informed framework incorporates orbital dynamics constraints into the
learning process, while the BiGRU-based architecture captures temporal
dependencies from both forward and backward directions for intention recognition.
The proposed method is validated through comprehensive numerical simulations.

\nocite{*}%% Remove this line from your manuscript.
\end{abstract}

% Use if graphical abstract is present
%\begin{graphicalabstract}
%\includegraphics{}
%\end{graphicalabstract}

% Research highlights
\begin{highlights}
\item Physics-informed deep learning framework for satellite trajectory prediction
\item BiGRU-based architecture for maneuver intention recognition
\item Integration of orbital dynamics constraints with data-driven models
\end{highlights}

% Keywords
% Each keyword is seperated by \sep
\begin{keywords}
Satellite cluster \sep Trajectory prediction \sep Intention recognition \sep
Physics-informed neural network \sep LSTM \sep BiGRU
\end{keywords}

\maketitle

% Main text

\section{Introduction}\label{sec:intro}

\section{Problem statement}\label{sec:problem}

\section{Physics-informed LSTM trajectory prediction}\label{sec:trajectory}

\section{BiGRU-based intention recognition}\label{sec:intention}

\section{Simulations and analysis}\label{sec:simulations}

\section{Conclusions}\label{sec:conclusions}

%% The Appendices part is started with the command \appendix;
%% appendix sections are then done as normal sections
%% \appendix

% To print the credit authorship contribution details
\printcredits

%% Loading bibliography style file
%\bibliographystyle{model1-num-names}
\bibliographystyle{cas-model2-names}

% Loading bibliography database
\bibliography{cas-refs}

% Biography
%\bio{}
% Here goes the biography details.
%\endbio

%\bio{pic1}
% Here goes the biography details.
%\endbio

\end{document}
